\documentclass[10pt,a4paper]{amsart}
\usepackage[margin=19mm]{geometry}
\usepackage{amsmath,amssymb,mathtools}
\usepackage[scr=rsfs,cal=euler]{mathalfa}
\usepackage{xfrac}
\usepackage[unicode,hidelinks,bookmarks=false]{hyperref}
\newcommand{\ie}{i.e.\ }
\newcommand{\cf}{cf.\ }
\newcommand{\resp}{respectively\ }
\newcommand{\Subsection}{\subsection}
\providecommand{\nobreakdash}{\nobreak\hbox{-}}
\setlength{\emergencystretch}{2em}
\multlinegap0pt
\usepackage{amssymb}
\usepackage{xcolor}
\usepackage{caption}
\newtheorem{proposition}{Proposition}
\newtheorem{theorem}{Theorem}
\newtheorem{claim}{Claim}
\title{Hanf Numbers for Poset Games}
\author{Fabi\'an Rivero Herrera}
\date{Source: arXiv:2504.07317v3 (CC BY 4.0)}
\begin{document}
\maketitle

\noindent\emph{Source and licence.} Hanf Numbers for Poset Games. Version arXiv:2504.07317v3 (CC BY 4.0), distributed by arXiv under the Creative Commons Attribution 4.0 International licence, http://creativecommons.org/licenses/by/4.0/, as recorded in the arXiv licence metadata for this exact version and shown on the abstract page https://arxiv.org/abs/2504.07317v3. Full licence notice: \texttt{licenses/arxiv-2504.07317-CC-BY-4.0.txt}.\par\smallskip

\noindent\textit{Editorial scope.} Counted unit: the article's theorem charwpog (source0.tex lines 225--227). The article's complete original proof of it is delivered with it (source0.tex lines 228--240, one \texttt{proof} environment of 13 lines). No mathematics is written, repaired or re-derived: the statement and the proof are the article's own lines, in the article's own notation, and no retained line is substituted or re-flowed.  Retained uncounted as the source-local context the proof needs: lines 143--145.  Nothing was omitted from any retained span, so the package carries no omission record.  No retained line contains a citation, so the wrapper carries no bibliography and no bibliography entry.  No retained line contains a figure environment, an image-inclusion command or drawing code, so no image is delivered and the package declares no asset dependency.  Editorial formatting only, in addition: an AMS \texttt{amsart} preamble that restates the article's own declarations and macros used by the retained lines, namely the environments \texttt{proposition}, \texttt{theorem}, \texttt{claim} on separate counters, together with the standard packages those lines need.  The retained environments print in this document as Proposition 1, Theorem 1, Claim 1 in order of appearance, whereas the article prints the same items with the numbering of its own full document.  The complete native source of the article as distributed in the arXiv e-print for this exact version is bundled unchanged as \texttt{sources/176281/source0.tex}.
\begin{proposition}\label{profini}
    If \((P; \leq_P)\) is a well partial order with a global minimum $0$, then, for any ordinal \(\sigma\), \((P^\sigma; \leq_{P^\sigma})\) is a well partial order and contains a global minimum.
\end{proposition}

\begin{theorem}\label{charwpog}
    Let $\Gamma$ be a finite set of ordinals containing a set $\{p_1,\dots,p_n\}$ of relatively prime natural numbers. If $(\mathcal{S}^{\overline{e}(\Gamma)+1};\leq_{\mathcal{S}^{\overline{e}(\Gamma)+1}})$ contains no infinite antichains, then $\Gamma$ is a $\mathsf{wpog}$.
\end{theorem}

\begin{proof}
    By transfinite induction on the exponent of the structures $(\mathcal{S}^{\sigma};\leq_{\mathcal{S}^{\sigma}})$. The base case is our hypothesis. Now, assume that $\Gamma\subset\omega^{\sigma}$ and that there are no infinite antichains in $(\mathcal{S}^{\sigma};\leq_{\mathcal{S}^{\sigma}})$. For any $\zeta \in \langle \Gamma \rangle^{\sigma + 1}$, either $\zeta\in\langle\Gamma\rangle^{\sigma}$ or there exist $\alpha \in \langle \Gamma \rangle^{\overline{e}(\Gamma) + 1}$ and $\xi \in \langle \Gamma \rangle^{\sigma}$ such that $\zeta = \omega^{\gamma} \otimes \alpha \oplus \xi$ for some $\gamma $ satisfying $\overline{e}(\omega^\gamma\otimes\alpha)=\sigma $. Let us consider an infinite set \( A= \{\alpha_i \otimes \omega^{\gamma} \oplus \xi_j \mid i,j \in \omega\}\subset\langle\Gamma\rangle^{\sigma+1}\), where $\overline{e}(\omega^\gamma\otimes\alpha_i)=\sigma $ and $\{\alpha_i \mid i\in\omega\}$ and $\{\xi_j \mid j\in\omega\}$ are arbitrary subsets of $\langle \Gamma \rangle^{\overline{e}(\Gamma)+1}$ and $\langle \Gamma \rangle^{\sigma}$, respectively. If for the terms with $\omega^\sigma$ of $A$ there are infinitely many distinct coefficients, then the set $\{\alpha_i \mid i\in\omega\}$ must be infinite. Hence we may find $i,i'\in\omega$ such that $\alpha' = \alpha_{i'} - \alpha_i \in \langle \Gamma \rangle^{\overline{e}(\Gamma)+1}$ and $\overline{e}(\alpha') = \overline{e}(\alpha_{i'})$. Consequently, \[ \alpha_{i'} \otimes \omega^{\gamma} \oplus \xi_{j'} -\alpha_i \otimes \omega^{\gamma} \oplus \xi_j = \alpha' \otimes \omega^{\gamma} \oplus \xi_{j'}, \] for arbitrary $j,j'$. Therefore, if there are infinitely many coefficients for the $\omega^\sigma$ term, the set $A$ cannot be an infinite antichain. Before proceeding, we will need the following claim.
\begin{claim}
    Let $\alpha=\omega^{\alpha_1} a_1 + \cdots + \omega^{\alpha_m} a_m\in \langle \Gamma \rangle^{\overline{e}(\Gamma)+1}$. For every ordinal $\beta$ satisfying $\alpha_m\leq\beta<\alpha_1 $, there exists $k_{\beta}\in \omega$ such that \[\alpha_\beta:=\omega^{\beta}k_{\beta}  + \omega^{\alpha_{i-1}} a_{i-1}+ \cdots + \omega^{\alpha_m} a_m \in\langle \Gamma \rangle^{\overline{e}(\Gamma)+1}, \] where $\alpha_{i-1}<\beta\leq\alpha_i$. If $\alpha_i-\alpha_{i-1}\geq\omega$ and $\alpha_{i-1}<\beta<\alpha_i$, we can take $k_\beta=0$.
\end{claim}
\begin{proof}[Proof of the claim] This is clear for $\alpha_m$. For $\alpha_m<\beta<\alpha_1$, consider the set $ \{\alpha\oplus\omega^{\beta}\otimes(b_1p_1+\cdots+b_np_n) \ |\ b_1,\dots,b_n\in\omega\}\subset\langle\Gamma \rangle^{\overline{e}(\Gamma)+1},$ which is not an antichain, by assumption. Therefore, we can find two comparable elements with difference $\omega^{\beta}d_\beta+\omega^{\alpha_{i-1}} a_{i-1} + \cdots + \omega^{\alpha_m} a_m\in\langle\Gamma \rangle^{\overline{e}(\Gamma)+1}$. Then, we take $k_{\beta} = d_\beta$.

For the second part, consider the set $\{\alpha \oplus \omega^{\alpha_{i-1}+k} p_1 \mid k \in \omega\}$. By similar arguments, we see that there are infinitely many elements of the form $\omega^{\alpha_{i-1}+k} p_1+ \omega^{\alpha_{i-1}} a_{i-1}+ \cdots + \omega^{\alpha_m} a_m $ in $\langle \Gamma \rangle^{\overline{e}(\Gamma)+1}$. Then, if we were unable to obtain $\omega^{\alpha_{i-1}} a_{i-1} + \cdots + \omega^{\alpha_m} a_m$, then there would have to be infinitely many generators in $\Gamma$, which is not the case, so $\omega^{\alpha_{i-1}} a_{i-1} + \cdots + \omega^{\alpha_m} a_m\in \langle \Gamma \rangle^{\overline{e}(\Gamma)+1}$.
\end{proof}
    From what we saw at the beginning, an antichain in $\langle\Gamma\rangle^{\sigma+1}$ would give us an antichain of the form $A=\{\alpha \otimes \omega^{\gamma} \oplus \xi_j \mid j \in \omega\}$ where $\Xi=\{\xi_j \mid j\in\omega\}$ is some subset of $\langle \Gamma \rangle^{\sigma } $ and $\alpha\in\langle \Gamma \rangle^{\overline{e}(\Gamma)+1}$ is fixed and satisfies $ \overline{e}(\omega^\gamma\otimes\alpha)=\sigma$. Assume that $\beta=\overline{e}(\Xi)\in\{\overline{e}(\xi_j)\mid j\in\omega\}$; that is,  for all $j\in\omega$, the first unfixed term of $\alpha \otimes \omega^{\gamma} \oplus \xi_j$ is of the form $\omega^\beta \cdot n_j^\beta$ whenever $\alpha_{i-1} < \beta < \alpha_i$, or $\omega^\beta \cdot (n_j^\beta + a_i)$ when $\beta = \alpha_i$. If there are infinitely many $n_j^\beta$, we can apply the claim and obtain $\{\alpha_\beta \otimes \omega^{\gamma} \oplus \xi_j \mid j \in \omega\}\subset\langle\Gamma\rangle^\sigma$. This cannot be an antichain, so we may compare two elements $\alpha_\beta \otimes \omega^{\gamma} \oplus \xi_j$ and $\alpha_\beta \otimes \omega^{\gamma} \oplus \xi_{j'}$. It is not hard to see that, consequently, we can also compare $\alpha \otimes \omega^{\gamma} \oplus \xi_j$ and $\alpha \otimes \omega^{\gamma} \oplus \xi_{j'}$, which is absurd. Therefore, we may assume that there are only finitely many distinct $n_j^\beta$’s and, in particular, we may fix one of them and obtain a new infinite antichain $A' \subseteq A$ whose first and second coefficients are fixed. If $\overline{e}(\Xi)>\overline{e}(\xi_j)$ for all $j\in\omega$, we apply the second part of the claim in a similar way.
    
    We may continue fixing coefficients, using the second part of the claim in the cases where $\alpha_i - \alpha_{i-1} \geq \omega$. After a finite number of steps, we will have fixed all coefficients between $\alpha_1$ and $\alpha_m$, and we will conclude that there is an infinite subset of $\{\xi_j \mid j \in \omega\}$ that forms an antichain, which is impossible. Hence, $A$ cannot be an antichain. Finally, if $\sigma$ is a limit ordinal, we proceed as in \hyperref[profini]{Proposition \ref{profini}}.
\end{proof}

\end{document}
